% 第8章 综合评价
\chapter{综合评价}

\begin{chapteroverview}
\textbf{这个分类是干什么的？}

综合评价解决的是"把多个指标合成一个总得分来排名"的问题。现实中评价一个对象往往要同时看好几个维度——评价门店要看坪效、毛利率、复购率、损耗率……每个指标量纲不同、方向不一，还可能有的"越大越好"有的"越小越好"。综合评价就是一套系统方法：先\textbf{给每个指标定权重}（谁更重要），再\textbf{把它们合成一个可比的总分}，最后排出优劣次序。

\textbf{美赛什么题型常用？}

综合评价是美赛\textbf{最高频}的建模题型之一。凡是题目出现"评价、排名、排序、优选、选最佳方案、确定指标权重"等字样，几乎都用这一章。其中\textbf{熵值法、AHP、TOPSIS}是三大核心方法，几乎每场比赛都至少用到一个；DEA用于效率评价，耦合协调度、障碍度模型用于发展质量诊断。

\textbf{学习路径建议}
先掌握\textbf{赋权}：客观赋权学\textbf{熵值法}（带星标，完全靠数据说话），主观赋权学\textbf{AHP层次分析法}（带星标，靠专家成对比较）；再学\textbf{CRITIC权重法}。然后学\textbf{TOPSIS}（带星标，与理想方案比距离排名）这个最常用的"打分器"，以及它和熵值法结合的\textbf{熵权TOPSIS}。\textbf{模糊综合评价}用于指标模糊、难以精确量化的场景。进阶部分了解DEA/SBM效率评价、BWM、耦合协调、障碍度、DEMATEL、ISM等专门工具。
\end{chapteroverview}

\section{基础级算法}

%% ========== 熵值法 ==========
\basicalgo{{\starmark}熵值法}{entropy_weight}

\begin{algointro}
熵值法是一种\textbf{完全客观}的定权方法：它不请专家打分，而是让数据自己说话——哪个指标在各评价对象之间差异越大、越能"拉开档次"，就给它越高的权重。最后按权重加权求和得到每个对象的综合得分并排名。
\end{algointro}

\begin{algoscene}
\begin{itemize}
    \item 美赛中需要\textbf{客观赋权}、避免主观偏好争议的综合评价
    \item 本案例：150家门店在坪效、毛利率、复购率、周转率、损耗率、投诉率六维下的经营质量排名
    \item 有较多样本（几十个以上）、各指标方向不一需要先同向化
    \item 论文中"客观权重确定"环节的标准方法
\end{itemize}
\end{algoscene}

\begin{algoprinciple}
"熵"本来是信息论里衡量\textbf{混乱程度}的量。熵值法的反直觉逻辑是：\textbf{一个指标在所有对象上取值都差不多，说明它区分不了谁好谁坏，信息量少、权重就低；反之取值越参差，越能拉开差距，权重就高}。

\begin{itemize}
    \item \textbf{第一步：方向同向化+归一化}。负向指标（如损耗率、投诉率，越低越好）先按"最大值平移"转成正向，再做min-max归一化到$[0,1]$。
    \item \textbf{第二步：算信息熵 $e_j$}。对第 $j$ 个指标，把各对象归一化取值转成比重 $p_{ij}$，算 $e_j=-k\sum p_{ij}\ln p_{ij}$。$e_j$ 越接近1，说明取值越均匀。
    \item \textbf{第三步：定权}。差异系数 $d_j=1-e_j$，再归一化 $w_j=d_j/\sum d_j$ 就是该指标权重。
    \item \textbf{第四步：合成得分}。$S_i=\sum_j w_j\cdot x'_{ij}$，把150家门店按 $S_i$ 降序排。
\end{itemize}

本案例权重：坪效0.284最高（门店间差距最大），其次库存周转率0.234、毛利率0.166；损耗率、投诉率因门店间差异小而权重低。
\end{algoprinciple}

\begin{algosposs}
\begin{enumerate}
    \item 在Dabbit AI SPSS工具箱中选择"熵值法"
    \item 上传评价数据：对象标识列选"门店编号"，纳入6个指标
    \item 设置参数：
    \begin{itemize}
        \item 指标方向：坪效/毛利率/会员复购率/库存周转率为正向；损耗率/顾客投诉率为负向
        \item 方向化方式：最大值平移；缺失处理整行删除
        \item 得分量纲：百分制换算；排序输出降序
        \item 常量指标权重记零并保留
    \end{itemize}
    \item 运行，查看权重表与综合得分排名
\end{enumerate}

\textbf{SPSS中实现}：SPSS无原生熵值法菜单，通过 \textbf{分析 → 描述统计} 配合变换计算，或用Dabbit AI SPSS工具箱一键完成。
\end{algosposs}

\begin{algoresult}
\begin{itemize}
    \item \textbf{权重表}：报告列出每个指标的熵值 $e_j$、差异系数 $d_j$、权重 $w_j$。本案例坪效权重0.284最高。
    \item \textbf{权重怎么解读}：权重高$\neq$指标业务上最重要，而是"它在样本间最能拉开差距"。这是熵值法的核心特点——客观但可能与直觉重要性不符。
    \item \textbf{综合得分排名}：本案例S115得分100分居首，S148得0分垫底，中间连续分布。
    \item \textbf{领先原因}：冠军店的优势主要落在权重最高（差异最大）的坪效等指标上。
\end{itemize}

\textbf{判断好坏}：权重之和=1，无指标权重为负；若某指标取值全相同（常量），其权重记0，应考虑剔除。
\end{algoresult}

\begin{algonotes}
\textbf{什么时候用？}需要客观、可复现、不掺杂专家主观的赋权排名；样本量较大。

\textbf{什么时候不用？}
\begin{itemize}
    \item 样本很少（几个对象）时，"差异"本身不稳定，权重不可靠
    \item 业务上某些指标虽差异小但极端重要（如安全性），纯数据赋权会低估它
\end{itemize}

\textbf{常见坑}
\begin{itemize}
    \item \textbf{坑1}：忘记先把负向指标同向化，直接拿"投诉率"算熵，方向就反了。
    \item \textbf{坑2}：归一化后出现0，导致 $\ln 0$ 报错——需加平移量。
    \item \textbf{坑3}：把"客观权重"当成"正确权重"。它反映区分度，不反映重要性，重要场合应与AHP等主客观方法结合。
\end{itemize}

\textbf{和相似算法的区别}：vs \textbf{AHP}——熵值法客观、靠数据，AHP主观、靠专家；vs \textbf{变异系数法/CRITIC}——都客观，但熵值法用信息熵，变异系数用标准差/均值，CRITIC兼顾对比强度与指标间冲突性。
\end{algonotes}

%% ========== AHP ==========
\basicalgo{{\starmark}层次分析法（AHP）}{ahp_method}

\begin{algointro}
AHP层次分析法是最经典的\textbf{主观赋权}方法。当指标间谁更重要难以直接量化时，它让专家只做"两两比较"（A和B哪个更重要、重要几倍），再把这些两两判断整理成一套权重，并检查专家判断是否自相矛盾。
\end{algointro}

\begin{algoscene}
\begin{itemize}
    \item 美赛中指标重要性无法从数据直接得出、需要专家经验判断时（选址、方案优选）
    \item 本案例：三家供应商在技术、价格、交付、服务、数据安全五准则下的择优
    \item 多准则决策、需要把定性经验转成定量权重
    \item 论文中"权重确定"环节，常与熵值法组合成主客观组合赋权
\end{itemize}
\end{algoscene}

\begin{algoprinciple}
让你直接给5个指标各打一个权重很难，但问你"技术和价格哪个更重要、重要多少倍"就容易多了。AHP就是把这些两两比较拼成一张判断矩阵，再反推出权重。

\begin{itemize}
    \item \textbf{构造判断矩阵}：用Saaty的1--9标度打分。$a_{ij}=1$ 表示两指标同等重要，$3$ 表示前者稍重要，$5$ 明显重要，$7$ 强烈重要，$9$ 极端重要；反之取倒数。
    \item \textbf{求权重}：对判断矩阵求最大特征根 $\lambda_{\max}$ 对应的特征向量，归一化即权重。
    \item \textbf{一致性检验（关键！）}：算一致性指标 $CI=(\lambda_{\max}-n)/(n-1)$，一致性比率 $CR=CI/RI$（$RI$ 是随机一致性指标，随阶数查表）。\textbf{$CR<0.1$ 才算专家判断前后一致、权重可用}；$CR\geq0.1$ 要回去改判断。
\end{itemize}

本案例准则层：技术能力权重0.317最高，价格0.260次之；综合后供应商A权重0.389最优。判断矩阵 $\lambda_{\max}=5.002$，$CR=0.001<0.1$，一致性很好。
\end{algoprinciple}

\begin{algosposs}
\begin{enumerate}
    \item 在Dabbit AI SPSS工具箱中选择"层次分析法AHP"
    \item 录入准则层两两判断（如"技术能力 vs 价格竞争力"的1--9标度）
    \item 若为双层结构，再逐准则录入"各方案 vs 该准则"的局部判断
    \item 设置参数：一致性检验阈值CR=0.10；专家判断聚合方式geometric（几何平均）
    \item 运行，查看权重向量、$\lambda_{\max}$、CI、CR一致性检验结果
\end{enumerate}

\textbf{SPSS中实现}：SPSS无原生AHP，通过Dabbit AI SPSS工具箱或R/Python扩展实现。
\end{algosposs}

\begin{algoresult}
\begin{itemize}
    \item \textbf{权重表}：每个准则的权重及排序。本案例技术能力0.317居首。
    \item \textbf{一致性比率CR}：本案例 $CR=0.001<0.1$，通过。若 $CR\geq0.1$，说明专家判断自相矛盾（如说A比B重要、B比C重要、C又比A重要），结果不可用。
    \item \textbf{综合权重}：方案综合权重 $=\sum$（准则权重 $\times$ 方案在该准则下局部权重）。本案例供应商A 0.389 $>$ B 0.319 $>$ C 0.292。
    \item \textbf{敏感性}：可对判断加小扰动看权重是否稳定；若两方案权重差小于0.03，视为排序接近。
\end{itemize}

\textbf{判断好坏}：CR$<$0.1是硬门槛；权重符合业务直觉、方案排序清晰即可采信。
\end{algoresult}

\begin{algonotes}
\textbf{什么时候用？}指标重要性主要靠经验判断、缺乏客观数据时；多方案比选。

\textbf{什么时候不用？}
\begin{itemize}
    \item 能拿到大量样本数据时，优先用熵值法等客观方法减少主观性
    \item 专家之间意见严重分歧时，AHP权重会不稳定
\end{itemize}

\textbf{常见坑}
\begin{itemize}
    \item \textbf{坑1}：不做一致性检验就报权重。CR超0.1的判断矩阵是无效的。
    \item \textbf{坑2}：指标过多（超过9个）时两两比较次数爆炸、一致性难保证。
    \item \textbf{坑3}：把AHP当客观方法。它本质是主观经验的量化，应说明专家构成。
\end{itemize}

\textbf{和相似算法的区别}：vs \textbf{熵值法}——AHP主观、熵值法客观，常组合赋权；vs \textbf{模糊层次分析FAHP}——AHP用精确数，FAHP用三角模糊数表达"大致重要"，更贴合专家模糊判断；vs \textbf{BWM}——BWM只需比较"最优/最劣"，比较次数更少、一致性更好。
\end{algonotes}

%% ========== CRITIC权重法 ==========
\basicalgo{CRITIC权重法}{critic_weight}

\begin{algointro}
CRITIC是一种比熵值法考虑更全面的客观赋权法：它不仅看指标自身的对比强度（标准差），还看指标之间的冲突性（相关程度）——既差异大、又与其他指标不重复的指标，权重最高。
\end{algointro}

\begin{algoscene}
\begin{itemize}
    \item 多个指标间可能信息重复，希望定权时兼顾"区分度"和"独立性"
    \item 本案例：150家供应商在报价、质量、交付、服务四维下的客观赋权
    \item 论文中想说明权重既考虑波动又考虑指标冗余
    \item 作为熵值法的改进版对照使用
\end{itemize}
\end{algoscene}

\begin{algoprinciple}
CRITIC考虑两个因素：

\begin{itemize}
    \item \textbf{对比强度}：指标标准差越大，对象间差异越大，信息量越大。
    \item \textbf{冲突性}：两个指标若高度相关（都在反映同一件事），信息重叠、贡献应打折；与其他指标相关越低，独立性越强、贡献越大。
    \item \textbf{信息量 $C_j$}$=\sigma_j\sum_k(1-r_{jk})$：标准差乘上"与其他指标不相关程度"之和。最后把 $C_j$ 归一化得到权重 $w_j$。
\end{itemize}

直觉：一个指标又能拉开差距、又不跟别人重复，才是真正"独当一面"的指标。
\end{algoprinciple}

\begin{algosposs}
\begin{enumerate}
    \item 在Dabbit AI SPSS工具箱中选择"CRITIC权重法"
    \item 上传数据，指定对象标识与各指标列、各指标方向
    \item 设置参数：标准化方式默认min-max；缺失值均值填补
    \item 运行，查看各指标标准差、相关阵、信息量 $C_j$ 与最终权重
\end{enumerate}
\end{algosposs}

\begin{algoresult}
\begin{itemize}
    \item \textbf{权重表}：列出每个指标的标准差、与其他指标的平均相关、信息量 $C_j$、权重 $w_j$。
    \item \textbf{解读}：标准差大且与其他指标相关低的指标权重最高；若某指标与别人高度相关，其权重被压下来。
    \item \textbf{后续}：把CRITIC权重代入TOPSIS或加权综合得分即可排名。
\end{itemize}
\end{algoresult}

\begin{algonotes}
\textbf{什么时候用？}指标间存在信息重叠、希望客观赋权同时去冗余。

\textbf{什么时候不用？}样本量小、相关阵估计不稳时慎用。

\textbf{常见坑}
\begin{itemize}
    \item \textbf{坑1}：不做指标方向同向化，负向指标混入会扭曲结果。
    \item \textbf{坑2}：混淆CRITIC与熵值法——前者兼顾相关性，后者只看离散度。
\end{itemize}

\textbf{和相似算法的区别}：vs \textbf{熵值法}——熵值法只用信息熵（离散度），CRITIC额外惩罚指标间相关；vs \textbf{独立性权系数法}——后者从复相关角度压冗余指标，思路相近。
\end{algonotes}

%% ========== 模糊综合评价 ==========
\basicalgo{模糊综合评价}{fuzzy_eval}

\begin{algointro}
很多评价是模糊的——"服务好不好"没法用一个精确数字衡量，只能说"很好/较好/一般/较差"。模糊综合评价用\textbf{隶属度}把这种"亦此亦彼"的模糊判断量化，再综合成一个评价等级。
\end{algointro}

\begin{algoscene}
\begin{itemize}
    \item 本案例：养老服务质量、满意度等难以精确打分的定性/半定性评价
    \item 指标评语是"优/良/中/差"等模糊等级
    \item 美赛中评价服务、生态、舒适度等带有主观模糊性的对象
    \item 需要给出"该对象综合属于哪一档"的结论
\end{itemize}
\end{algoscene}

\begin{algoprinciple}
想象让你给"服务水平"打分，你很难说它到底是8.3还是8.7分，但你可以说：\textbf{它"隶属"于"优"的程度是0.2、"良"0.5、"中"0.3}——这就是隶属度。

\begin{itemize}
    \item \textbf{第一步：定因素集和评语集}。因素集如\{可及性、专业性、规范性……\}，评语集如\{优、良、中、差\}。
    \item \textbf{第二步：建隶属度矩阵 $R$}。对每个因素，专家/数据给出它对各评语的隶属度（一行和为1）。
    \item \textbf{第三步：定权重向量 $W$}。可用AHP或熵值法得到各因素权重。
    \item \textbf{第四步：模糊合成}。$B=W\circ R$（矩阵乘法或模糊算子），得到对象对各评语的综合隶属度。
    \item \textbf{第五步：最大隶属度原则}。$B$ 中最大值对应的评语，就是综合评价等级。
\end{itemize}
\end{algoprinciple}

\begin{algosposs}
\begin{enumerate}
    \item 在Dabbit AI SPSS工具箱中选择"模糊综合评价"
    \item 录入评价因素、评语等级，以及各因素对各等级的隶属度
    \item 设定因素权重（可对接AHP/熵值法结果）
    \item 运行，查看综合隶属度向量与最终判定等级
\end{enumerate}
\end{algosposs}

\begin{algoresult}
\begin{itemize}
    \item \textbf{综合隶属度向量 $B$}：如 (0.2, 0.5, 0.2, 0.1)，分别对应优/良/中/差。
    \item \textbf{判定等级}：取最大分量对应等级——本例最大0.5在"良"，故综合评为"良"。
    \item \textbf{注意}：若最大两个分量很接近（如0.45 vs 0.42），说明对象处于两档之间，不要过度解读。
\end{itemize}
\end{algoresult}

\begin{algonotes}
\textbf{什么时候用？}评价标准模糊、定性指标多、需要给出等级判定时。

\textbf{什么时候不用？}指标全部精确可测、直接算分更清晰时，不必套模糊。

\textbf{常见坑}
\begin{itemize}
    \item \textbf{坑1}：隶属度矩阵每一行不归一（和不为1），合成结果失真。
    \item \textbf{坑2}：最大隶属度过低（如只占0.3）却硬判一档，应说明"接近临界"。
\end{itemize}

\textbf{和相似算法的区别}：vs \textbf{AHP}——AHP解决"权重"，模糊综合解决"模糊评价本身"，二者常组合（FAHP是把AHP模糊化）。
\end{algonotes}

%% ========== TOPSIS ==========
\basicalgo{{\starmark}TOPSIS法}{topsis_method}

\begin{algointro}
TOPSIS（逼近理想解排序法）是综合评价里\textbf{最常用的打分器}。它的思路很直观：在所有方案中找一个"完美方案（正理想解，每项都最好）"和一个"最差方案（负理想解，每项都最差）"，然后看每个对象\textbf{离谁更近}——离正理想越近、离负理想越远，就越好。
\end{algointro}

\begin{algoscene}
\begin{itemize}
    \item 美赛中\textbf{排名、优选、方案排序}的头号方法，几乎万能
    \item 本案例：150家门店在六项经营指标下的综合排序
    \item 多指标、多对象，需要一个连续可比的综合得分
    \item 常与熵值法组合成"熵权TOPSIS"，是论文高频组合
\end{itemize}
\end{algoscene}

\begin{algoprinciple}
想象一场选秀：评委心里有个"全能冠军模板"（每项都是全场最佳）和一个"全能垫底模板"（每项都最差）。每个选手看他\textbf{离冠军模板近不近、离垫底模板远不远}。

\begin{itemize}
    \item \textbf{第一步：向量归一化+加权}。把各指标归一化，乘上权重（等权或熵权）。
    \item \textbf{第二步：确定正负理想解}。正理想解 $Z^+$ = 各指标最优值，负理想解 $Z^-$ = 各指标最差值（先统一方向，负向指标越小越好）。
    \item \textbf{第三步：算距离}。每个对象到正理想解的欧氏距离 $D^+$，到负理想解的距离 $D^-$。
    \item \textbf{第四步：算贴近度} $C_i=\dfrac{D_i^-}{D_i^+ + D_i^-}$。$C_i$ 越接近1越接近正理想、越好；越接近0越差。
\end{itemize}

本案例在加权归一化空间中，对象M0067贴近度1.0000（即它本身就是正理想解）排名第一。
\end{algoprinciple}

\begin{algosposs}
\begin{enumerate}
    \item 在Dabbit AI SPSS工具箱中选择"TOPSIS法"
    \item 上传数据，指定对象标识列与各指标列
    \item 设置参数：
    \begin{itemize}
        \item 指标方向：正向（越大越好）/负向（越小越好，先同向化）
        \item 权重：等权，或接入熵权/CRITIC权重
        \item 归一化方式：向量归一化
    \end{itemize}
    \item 运行，查看正/负理想解、各对象距离 $D^+,D^-$ 与贴近度 $C$ 排名
\end{enumerate}

\textbf{SPSS中实现}：SPSS无原生TOPSIS，通过Dabbit AI SPSS工具箱或R/Python扩展实现。
\end{algosposs}

\begin{algoresult}
\begin{itemize}
    \item \textbf{正/负理想解}：列出各指标的最优、最差值。
    \item \textbf{$D^+,D^-$}：每个对象到两个理想解的距离。
    \item \textbf{贴近度 $C$}：核心结果，$C\in[0,1]$，按 $C$ 降序排名。本案例M0067的 $C=1.0000$（与正理想重合）居首。
    \item \textbf{稳健性}：可做留一敏感性（去掉某指标重排），看名次是否稳定。
\end{itemize}

\textbf{判断好坏}：$C$ 是相对值，只在本样本内可比；$C=1$ 即该对象在所有指标上都最优。
\end{algoresult}

\begin{algonotes}
\textbf{什么时候用？}多对象多指标的排名优选，是综合评价的"通用底座"。

\textbf{什么时候不用？}
\begin{itemize}
    \item 指标间存在不可公度、或有缺失严重时，先清洗
    \item 需要判断"属于哪一档"而非连续排名时，模糊综合或RSR更合适
\end{itemize}

\textbf{常见坑}
\begin{itemize}
    \item \textbf{坑1}：负向指标不同向化就直接算距离，方向错误。
    \item \textbf{坑2}：用等权却声称"客观"。等权也是一种主观假设，建议配合熵权。
    \item \textbf{坑3}：只看贴近度绝对值。$C$ 是相对排名工具，不同批次样本的 $C$ 不能直接比大小。
\end{itemize}

\textbf{和相似算法的区别}：vs \textbf{熵权TOPSIS}——本TOPSIS权重可等权可主观，熵权TOPSIS用熵值法客观定权；vs \textbf{VIKOR}——TOPSIS看距离理想解的贴近度，VIKOR是折中排序，引入群体效用与个体遗憾。
\end{algonotes}

%% ========== 熵权TOPSIS ==========
\basicalgo{熵权TOPSIS}{entropy_topsis}

\begin{algointro}
熵权TOPSIS是\textbf{熵值法定权 + TOPSIS排序}的强强组合：先用熵值法从数据客观算出各指标权重，再把这套权重喂给TOPSIS算贴近度排名。它兼顾了客观赋权和距离排序，是美赛综合评价的"黄金组合"。
\end{algointro}

\begin{algoscene}
\begin{itemize}
    \item 美赛中最常被推荐的综合评价"标准答案"之一
    \item 本案例：150家连锁门店在销售额、坪效、复购率、满意度、成本、损耗六维下的运营绩效排名
    \item 既不想主观定权、又想要连续可比的贴近度得分
    \item 论文中权重客观、排序清晰、可做稳健性检验
\end{itemize}
\end{algoscene}

\begin{algoprinciple}
两步走：

\begin{itemize}
    \item \textbf{第一步：熵法定权}。对各指标算信息熵 $e_j$、效用值 $d_j=1-e_j$、熵权 $w_j=d_j/\sum d$。某指标对象间差异越大、熵越小、权重越高。本案例商品损耗率权重0.275最高、顾客满意度0.005最低。
    \item \textbf{第二步：TOPSIS排序}。把归一化矩阵乘熵权，找正/负理想解，算 $D^+,D^-$，贴近度 $C=D^-/(D^++D^-)$。
\end{itemize}

负向指标（运营成本、商品损耗率）先按最大值差法 $max-x$ 同向化。本案例M0067贴近度1.0000居首，M0100次之0.810。
\end{algoprinciple}

\begin{algosposs}
\begin{enumerate}
    \item 在Dabbit AI SPSS工具箱中选择"熵权TOPSIS"
    \item 上传数据，指定对象标识列（门店编号）
    \item 设置参数：
    \begin{itemize}
        \item 方向映射：商品损耗率、运营成本为负向（negative），其余正向
        \item 负向变换：max\_minus；归一化：vector向量归一化
        \item 缺失处理：median\_impute中位数填补
        \item 敏感性：leave\_one\_out留一法
    \end{itemize}
    \item 运行，查看熵权表、贴近度排名与留一稳健性
\end{enumerate}
\end{algosposs}

\begin{algoresult}
\begin{itemize}
    \item \textbf{熵权表}：本案例商品损耗率0.275、运营成本0.243权重最高（最能区分门店），顾客满意度0.005最低。
    \item \textbf{贴近度排名}：M0067为1.0000、M0100为0.8100……按 $C$ 降序。
    \item \textbf{留一稳健性}：逐一去掉指标重排，若头部名次稳定则结论可信。
\end{itemize}

\textbf{判断好坏}：权重反映区分度（非重要性）；排名对留一指标不敏感即稳健。
\end{algoresult}

\begin{algonotes}
\textbf{什么时候用？}希望权重完全客观、又要连续贴近度排名时的首选。

\textbf{什么时候不用？}样本很少、或业务上必须体现某些指标重要性时，纯熵权会失真。

\textbf{常见坑}
\begin{itemize}
    \item \textbf{坑1}：忽略负向指标同向化，成本/损耗方向算反。
    \item \textbf{坑2}：把熵权当重要性权重。它只代表区分度，低权重的满意度可能业务上极重要。
\end{itemize}

\textbf{和相似算法的区别}：它是熵值法与TOPSIS的串联；若想权重更贴合业务，可换成AHP权重再代入TOPSIS。
\end{algonotes}

\section{进阶级算法速览}

%% ========== SBM ==========
\advancedalgo{SBM模型}{sbm_model}

\begin{algobrief}
\textbf{一句话}：考虑了"松弛变量"和"非期望产出"的DEA效率评价模型，能同时处理少投入、多产出、少排污。

\textbf{美赛场景区}：绿色生产效率评价，如制造业企业在投入人力能耗、产出产值的同时排放二氧化碳。

\textbf{核心原理}：用Tone的非径向SBM模型，把投入冗余、期望产出不足、非期望产出过多的松弛量直接计入目标函数，经线性规划逐单元求效率。本案例150家企业平均效率0.731，21家有效，冗余主要在碳排放（松弛率26.2\%）。

\textbf{操作要点}：工具箱选"SBM"，指定投入、期望产出、非期望产出，VRS，非期望产出弱可处置。

\textbf{结果关键}：效率得分0--1，越接近1越好；看松弛率定位冗余来源（投入侧还是排放侧）。

\textbf{注意}：效率是组内相对值，跨期跨组不可直接比；需线性规划求解，SPSS原生不支持，用R/Python扩展。
\end{algobrief}

%% ========== 变异系数法 ==========
\advancedalgo{变异系数法（信息量权重法）}{cv_weight}

\begin{algobrief}
\textbf{一句话}：用"标准差/均值"（变异系数）衡量指标在样本间的相对波动，波动越大权重越高的客观赋权法。

\textbf{美赛场景区}：完全不引入主观打分，依据数据自身区分信息为门店经营指标定权。

\textbf{核心原理}：逐列算均值与标准差，变异系数$=$标准差/均值，归一化为权重，再对极差归一化数据加权求和。本案例"线上销售占比"权重最高（0.2805），综合得分最高对象M1142（0.7068）。

\textbf{操作要点}：工具箱选"变异系数法"，指标列自动识别，负向指标倒数变换，缺失均值填补。

\textbf{结果关键}：变异系数大的指标权重高；注意要求指标取值为正才能算CV。

\textbf{注意}：与熵值法同为客观赋权，思路相近但CV用离散系数、熵值用信息熵；量纲不同时要先同向化。
\end{algobrief}

%% ========== 模糊层次分析法 ==========
\advancedalgo{模糊层次分析法}{fahp_method}

\begin{algobrief}
\textbf{一句话}：把AHP中专家的"点估计"判断升级为"三角模糊数(l,m,u)"，表达"大概在这区间"的模糊判断，更贴近真实决策心理。

\textbf{美赛场景区}：多准则供应商择优，专家难以用精确数比较、判断带模糊性时。

\textbf{核心原理}：成对比较表示为三角模糊数$(l,m,u)$，用中值矩阵做一致性检验（CR$<$0.1），再用Buckley几何平均法算模糊权重、去模糊化归一化排序，并做对数扰动敏感性。本案例供应商A综合权重0.389最优，技术能力准则权重0.317。

\textbf{操作要点}：工具箱选"模糊层次分析"，去模糊化centroid，权重geometric\_mean，CR阈值0.10，专家几何聚合。

\textbf{结果关键}：CR$<$0.1通过；报告模糊权重区间与去模糊权重，权重差小于0.03视为排序接近。

\textbf{注意}：它是AHP的模糊升级版，适合判断不确定的场景；计算更复杂，需R/Python扩展。
\end{algobrief}

%% ========== RFM ==========
\advancedalgo{RFM客户价值模型}{rfm_model}

\begin{algobrief}
\textbf{一句话}：用最近消费时间(R)、消费频率(F)、消费金额(M)三个维度给客户打分分层，识别高价值客户。

\textbf{美赛场景区}：电商/零售客户分群与差异化营销，如找出贡献大部分营收的核心会员。

\textbf{核心原理}：从交易流水按客户聚合出R/F/M，分位数各打1--5分，以每维中位数为界分高低，组合成八类客群（重要价值/发展/保持/挽留等）。本案例重要价值客户占29.8\%却贡献68.3\%金额。

\textbf{操作要点}：工具箱选"RFM"，指定客户编号、消费日期、金额列，参考日期取数据最近日，等权，中位数分界。

\textbf{结果关键}：各客群的数量占比、金额占比、价值强度；帕累托曲线看集中度。

\textbf{注意}：R越小越好（越近越活跃）、F和M越大越好；权重默认等权，重要业务可调整。
\end{algobrief}

%% ========== VIKOR ==========
\advancedalgo{VIKOR法}{vikor_method}

\begin{algobrief}
\textbf{一句话}：一种"折中"排序法，在群体效用最优和个体遗憾最小之间权衡，给出妥协最优解。

\textbf{美赛场景区}：多指标供应商/方案选优，指标相互冲突、需要可接受的折中方案时。

\textbf{核心原理}：把各方案到理想点的加权距离拆成群体效用$S$和个体遗憾$R$，以折中系数$v=0.5$合成$Q$，再用可接受优势和稳定性两个条件裁决唯一最优。本案例S090的$Q=0$为唯一折中最优。

\textbf{操作要点}：工具箱选"VIKOR"，等权或自定义权重，极差归一化，$v=0.5$，方向自动识别。

\textbf{结果关键}：$Q$越小越好；需同时通过可接受优势（$Q$差达阈值）与稳定性（$S$或$R$也最优）检验。

\textbf{注意}：与TOPSIS同属多属性决策，但VIKOR显式引入"遗憾"和妥协条件，适合冲突型方案选择。
\end{algobrief}

%% ========== 功效系数法 ==========
\advancedalgo{功效系数法}{efficacy_coef}

\begin{algobrief}
\textbf{一句话}：把每个指标按"不允许值→满意值"线性映射到60--100分，再加权汇总成百分制综合得分。

\textbf{美赛场景区}：门店服务质量、政府绩效等需要百分制打分和分等级（优/良/中/及格）的评价。

\textbf{核心原理}：为每个指标定满意值（95分位）和不允许值（5分位），单项功效系数$=60+40\times$(实际-不允许)/(满意-不允许)，越界截断，加权求和。本案例门店056得分98.47（优），门店140为62.94（及格）。

\textbf{操作要点}：工具箱选"功效系数法"，指定正向/逆向指标，满意/不允许分位点0.95/0.05，范围60--100，越界截断。

\textbf{结果关键}：综合得分+等级+短板指标；平均功效系数低的指标是共性短板。

\textbf{注意}：基准值（满意/不允许）设定主观，要说明依据；它输出百分制便于管理沟通。
\end{algobrief}

%% ========== BWM ==========
\advancedalgo{最好最差法（BWM）}{bwm_method}

\begin{algobrief}
\textbf{一句话}：只需专家指定"最好"和"最差"准则，分别比较它们与其他准则，比较次数比AHP少、一致性更好的主观赋权法。

\textbf{美赛场景区}：多专家参与的指标权重确定，如30人评审组确定供应商评价准则权重。

\textbf{核心原理}：专家给出最优准则到各准则的BO向量、各准则到最劣准则的OW向量，用最小化最大偏差的线性规划求权重，以$CR=\xi^*/CI$检验一致性。本案例产品质量权重0.421最高，30组判断一致性均值0.0083全部通过。

\textbf{操作要点}：工具箱选"BWM"，指定best/worst准则，几何聚合，一致性阈值0.10。

\textbf{结果关键}：$CR<0.1$通过；权重和为1；报告专家间分歧幅度。

\textbf{注意}：是AHP的改进，比较次数从$n(n-1)/2$降到$2n-3$；仍属主观赋权，需R/Python扩展。
\end{algobrief}

%% ========== 障碍度模型 ==========
\advancedalgo{障碍度模型}{obstacle_model}

\begin{algobrief}
\textbf{一句话}：不只算综合得分，更专门诊断"到底是哪些指标在拖后腿、阻碍系统提升"。

\textbf{美赛场景区}：县域发展、城市宜居等综合评价的短板诊断，回答"该优先抓哪个指标/子系统"。

\textbf{核心原理}：极差标准化后用"1-标准化值"表示偏离最优的差距，障碍度$=$偏离度$\times$权重占全部加权偏离的比例；累计占比首次达60\%的指标为主要障碍因子。本案例"乡村旅游收入占比"障碍度最高（17.12\%），产业融合子系统合计35.4\%。

\textbf{操作要点}：工具箱选"障碍度模型"，指定对象列、指标列、准则层归属，权重entropy，缺失均值填补。

\textbf{结果关键}：按平均障碍度降序找主要障碍因子；子系统障碍度找短板模块。

\textbf{注意}：它与综合得分互补——得分告诉你"谁差"，障碍度告诉你"差在哪、先改哪"。
\end{algobrief}

%% ========== 德尔菲法 ==========
\advancedalgo{德尔菲法（Delphi）}{delphi_method}

\begin{algobrief}
\textbf{一句话}：通过多轮匿名征求专家意见、逐轮反馈统计结果，直到意见收敛，从而筛选出高共识的评价条目。

\textbf{美赛场景区}：构建评价指标体系之初，筛选哪些指标真正重要、专家意见一致。

\textbf{核心原理}：多轮匿名打分（1--5分），每轮算均值、满分频率、变异系数、四分位距，用Kendall协调系数$W$衡量专家整体协调，按均值和变异系数双阈值筛选条目，直至收敛。本案例3轮后$W=0.874$，保留3/5条目。

\textbf{操作要点}：工具箱选"德尔菲法"，录入多轮专家评分，均值阈值3.5、CV阈值0.25，Kendall $W$检验。

\textbf{结果关键}：末轮$W$接近1且显著、均值位移小说明收敛；保留达标的条目。

\textbf{注意}：它是指标\textbf{筛选}工具，本身不算综合得分；常作为AHP/模糊综合的前置步骤。
\end{algobrief}

%% ========== 灰色关联分析 ==========
\advancedalgo{灰色关联分析}{grey_relational}

\begin{algobrief}
\textbf{一句话}：以一个参考序列为基准，算各比较序列与它的"形状相似度"，判断哪些因素与结果最同步。

\textbf{美赛场景区}：小样本下识别哪些经营因素与销售额走势最贴近、谁是主导驱动。

\textbf{核心原理}：先极差标准化，指定参考列（如日销售额），逐点算绝对差和关联系数，加权平均成灰色关联度并排序，分辨系数$\rho=0.5$。本案例日客流量关联度最高（$r=0.917$）。

\textbf{操作要点}：工具箱选"灰色关联分析"，指定参考列，分辨系数0.5，等权，极差标准化。

\textbf{结果关键}：关联度越大越贴近参考；$\rho$在0.1--0.9变化排序稳定即可靠。

\textbf{注意}：它衡量"形状同步"而非因果，适合样本少、信息不完全的小样本系统分析。
\end{algobrief}

%% ========== 独立性权系数法 ==========
\advancedalgo{独立性权系数法}{independence_weight}

\begin{algobrief}
\textbf{一句话}：认为与其他指标重叠（相关）越多的指标冗余越重、权重应越小，按信息独立性客观定权。

\textbf{美赛场景区}：指标间高度相关（如客流、销售额、坪效互相牵连）时，给门店评价定一组不重复信息的权重。

\textbf{核心原理}：算指标相关阵，以复相关系数衡量该指标被其余指标线性解释的程度，按$1/$复相关归一化得权重。本案例顾客满意度复相关最低（0.766）获最高权重18.6\%。

\textbf{操作要点}：工具箱选"独立性权系数法"，标准化minmax，pearson相关，复相关口径。

\textbf{结果关键}：复相关越低越独立、权重越高；权重反映信息独立性而非业务重要性。

\textbf{注意}：与CRITIC思路相近（都惩罚相关），但独立性权直接用复相关，CRITIC用标准差$\times$(1-相关)。
\end{algobrief}

%% ========== ISM ==========
\advancedalgo{解释结构模型（ISM）}{ism_model}

\begin{algobrief}
\textbf{一句话}：把一堆相互影响的因素，通过邻接矩阵和可达矩阵，梳理成"表层结果—中层传导—深层根源"的层级结构。

\textbf{美赛场景区}：分析影响工程质量/系统性能的众多因素，找出哪些是根源、哪些是表象。

\textbf{核心原理}：由专家两两判断建0/1邻接矩阵，加单位阵做布尔传递闭包得可达矩阵，用"可达集$\subseteq$先行集"逐层剥离分层，再做MICMAC驱动-依赖分类。本案例20因素分5层，高层领导重视、制度建设是根源驱动。

\textbf{操作要点}：工具箱选"ISM"，录入因素间二元关系，一致率阈值0.5，含MICMAC分类。

\textbf{结果关键}：层级图（表层L1到根源L5）；驱动型因素优先治理，依赖型只是结果。

\textbf{注意}：它梳理因素\textbf{结构关系}，不打分排名；判断矩阵来自专家，需多人共识。
\end{algobrief}

%% ========== 耦合协调度模型 ==========
\advancedalgo{耦合协调度模型}{coupling_coord}

\begin{algobrief}
\textbf{一句话}：衡量两个（或多个）子系统之间"相互配合、协调发展"的程度，输出协调等级和短板系统。

\textbf{美赛场景区}：如城镇化与生态环境、经济与社会两系统的协同发展评估，面板时序数据。

\textbf{核心原理}：子系统内用熵权合成综合指数$U$，算耦合度$C$、综合发展指数$T$、耦合协调度$D=\sqrt{C\times T}$，按0.1间隔分十级（失调—协调），按指数差判断滞后系统。本案例$D$均值0.693（中级协调），逐期+0.014。

\textbf{操作要点}：工具箱选"耦合协调度"，指定时间/个体列、各子系统指标组、负向指标，功效系数标准化+熵权。

\textbf{结果关键}：$D$的协调等级、随时间趋势斜率、哪个子系统滞后。

\textbf{注意}：耦合度高不代表发展好（低水平也可耦合），必须结合$T$看$D$；需R/Python扩展。
\end{algobrief}

%% ========== Malmquist指数 ==========
\advancedalgo{Malmquist指数}{malmquist_index}

\begin{algobrief}
\textbf{一句话}：基于DEA前沿，衡量决策单元全要素生产率的跨期变化，并拆成"效率追赶"和"技术进步"两部分。

\textbf{美赛场景区}：评价多家门店/企业连续几年生产率是进步还是退步、靠管理追赶还是技术升级。

\textbf{核心原理}：用DEA求各期前沿距离函数，按Färe分解算Malmquist指数$M=$技术效率变化$\times$技术进步，再细分纯技术效率与规模效率变化。本案例平均每期提升6.5\%，主要靠技术进步（1.089）。

\textbf{操作要点}：工具箱选"Malmquist"，指定DMU、时期、投入列、产出列，输入导向，几何平均汇总。

\textbf{结果关键}：$M>1$进步、$<1$退步；分解看是EFFCH还是TECHCH主导。

\textbf{注意}：需平衡投入产出面板；它是DEA的动态版，静态效率用DEA/SBM。
\end{algobrief}

%% ========== DEA ==========
\advancedalgo{数据包络分析（DEA）}{dea_method}

\begin{algobrief}
\textbf{一句话}：不用预设权重和生产函数，靠线性规划构造"生产前沿"，评价多投入多产出决策单元的相对效率。

\textbf{美赛场景区}：门店/银行/医院等单位投入资源转化为产出的效率评价，找标杆和低效改进空间。

\textbf{核心原理}：用包络线性规划构造由最有效单元围成的前沿面，测各单元到前沿的径向距离得效率分（0--1），再分解纯技术效率与规模效率。本案例150店中34家前沿有效，平均效率0.888，低效店平均有14.4\%投入压缩空间。

\textbf{操作要点}：工具箱选"DEA"，指定DMU列、投入列、产出列，投入导向、VRS，样本量$\geq 3(m+s)$。

\textbf{结果关键}：效率=1为有效；报告松弛量、规模报酬类型与参照标杆单元。

\textbf{注意}：效率是组内相对值，跨期不可直接比；不区分随机噪声（SBM/SFA改进）；需R/Python扩展。
\end{algobrief}

%% ========== 综合指数 ==========
\advancedalgo{综合指数}{composite_index}

\begin{algobrief}
\textbf{一句话}：把多个量纲方向不一的指标同向化、标准化后加权平均成一个0--1的综合指数，简单直接地排名分档。

\textbf{美赛场景区}：多区域门店经营绩效的统一打分、五档分档和区域比较。

\textbf{核心原理}：标定各指标方向（正向/逆向），统一变换后min-max标准化，按权重加权平均得0--1指数，降序排名并按分位数分五档。本案例ST0056指数0.959最高。

\textbf{操作要点}：工具箱选"综合指数"，指定id列、指标列与方向，加权平均，等权，分位数分档，minmax。

\textbf{结果关键}：综合指数均值/区间、五档分布、领先与落后对象。

\textbf{注意}：这是最朴素的加权综合，权重主观；追求客观可换熵权，追求距离可比可用TOPSIS。
\end{algobrief}

%% ========== DEMATEL ==========
\advancedalgo{DEMATEL}{dematel_method}

\begin{algobrief}
\textbf{一句话}：从专家两两影响打分中，算每个因素的"影响度、被影响度、中心度、原因度"，区分原因因素和结果因素。

\textbf{美赛场景区}：供应链/系统改进中找"牵一发而动全身"的关键源头因素。

\textbf{核心原理}：由专家0--4打分建直接影响矩阵，规范化后求综合影响矩阵，中心度$M=D+R$（谁关键）、原因度$N=D-R$（正为原因、负为结果）。本案例信息共享水平中心度最高，是核心原因因素。

\textbf{操作要点}：工具箱选"DEMATEL"，0--4级直接打分，最大行和归一化，阈值取矩阵均值。

\textbf{结果关键}：中心度排序找关键因素；原因度正负区分原因/结果，优先改原因因素。

\textbf{注意}：原因/结果是打分相对属性，不等于经过检验的因果链；与ISM互补——ISM分层级、DEMATEL判因果属性。
\end{algobrief}

%% ========== RSR ==========
\advancedalgo{秩和比法（RSR）}{rsr_method}

\begin{algobrief}
\textbf{一句话}：把各指标先编秩（排名次），再把秩次加权平均成一个0--1的综合秩和比，最后分档。

\textbf{美赛场景区}：医院科室、多指标评价的综合排序与相对档位划分，尤其适合指标量纲差异大时。

\textbf{核心原理}：高优指标从小到大编秩、低优指标反向，秩次按权重相加并归一得RSR，用概率单位回归检验其近似正态，再按分档概率划档。本案例科室011的RSR=0.989最高，分较差/中/良/优四档。

\textbf{操作要点}：工具箱选"秩和比法"，指定对象列、低优列，并列取平均秩，等权，正态性检验开启。

\textbf{结果关键}：RSR越大越好；probit回归$R^2$高说明近似正态、分档可信。

\textbf{注意}：编秩会损失原始数值信息（只保留顺序）；它简单稳健，适合数据质量一般的综合评价。
\end{algobrief}
